%===========================================
% TEMPLATE LaTeX ABNT - CLASSE MEMOIR
% Compilável e modular, sem gambiarras
%===========================================
\documentclass[12pt,a4paper,oneside]{memoir}
\usepackage[brazilian]{babel}
\usepackage{fontspec}
\usepackage{tikz}

\setmainfont{CMU Serif}
\setlrmarginsandblock{3cm}{2cm}{*}
\setulmarginsandblock{3cm}{2cm}{*}
\checkandfixthelayout
\OnehalfSpacing
\setlength{\parindent}{1.25cm}

%===========================================
% ESTILO DE CAPÍTULO NUMERADO (Com TikZ)
%===========================================
\makechapterstyle{CapitulosNumerado}{
    \setlength{\beforechapskip}{0.46cm}
    \renewcommand{\printchaptername}{}
    \renewcommand{\chapternamenum}{}
    \renewcommand{\afterchapternum}{\par\nobreak\vskip 1.3cm}
    \renewcommand{\printchapternum}{%
        \noindent
        \begin{tikzpicture}[remember picture, overlay]
            \def\numsize{61}
            \def\barwidth{0.8pt}
            \def\leftoffset{2.3cm}

            \node[inner sep=-1pt, anchor=south west, opacity=0] (num) at (\leftoffset, 0)
             {\fontsize{\numsize}{\numsize}\selectfont\rmfamily\thechapter};
            \node[inner sep=0pt, anchor=south west] at (\leftoffset, -1pt)
             {\fontsize{\numsize}{\numsize}\selectfont\rmfamily\thechapter};
            \node[inner sep=0pt, anchor=south west, font=\Large\rmfamily] (chaptext)
                 at (0pt, -4pt) {\chaptername};

            \coordinate (StartPoint) at (0.5*\barwidth, 13pt);
            \draw[line width=\barwidth] (StartPoint) -- (StartPoint |- num.north) -- (num.north west);
            \coordinate (RightWall) at (\linewidth, 0);
            \draw[line width=\barwidth] (num.north east) -- (RightWall |- num.north) -- (RightWall |- num.south) -- (num.south east);
        \end{tikzpicture}%
    }
    \renewcommand{\printchaptertitle}[1]{%
        \raggedright\huge\bfseries ##1
    }
}

%===========================================
% ESTILO DE CAPÍTULO NÃO NUMERADO (Centralizado com linha)
%===========================================
\makechapterstyle{estiloNaoNumerado}{
    \setlength{\beforechapskip}{0pt}
    \setlength{\afterchapskip}{2cm}
    \renewcommand{\printchaptername}{}
    \renewcommand{\chapternamenum}{}
    \renewcommand{\printchapternum}{}
    \renewcommand{\printchaptertitle}[1]{%
        \begin{center}
            \huge\bfseries\MakeUppercase{##1}\\[0.5cm]
            \rule{\textwidth}{0.3pt}
        \end{center}
    }
}

%===========================================
% COMANDO MODULAR: Capítulo Centralizado
%===========================================
\newcommand{\capituloCentralizado}[1]{%
    \chapterstyle{estiloNaoNumerado}%
    \chapter*{#1}%
    \addcontentsline{toc}{chapter}{\MakeUppercase{#1}}%
    \chapterstyle{CapitulosNumerado}%
}

%===========================================
% COMANDO PARA TÍTULOS DIVISÓRIOS (sem número de página no TOC)
%===========================================
\newcommand{\tituloDivisorio}[1]{%
    \chapterstyle{estiloNaoNumerado}%
    \chapter*{#1}%
    \addtocontents{toc}{\protect\vspace{0.5cm}}%
    \addtocontents{toc}{\protect\mbox{}\protect\hfill\textbf{\MakeUppercase{#1}}\protect\hfill\protect\mbox{}\par}%
    \addtocontents{toc}{\protect\vspace{0.3cm}}%
    \chapterstyle{CapitulosNumerado}%
}

%===========================================
% MACROS PARA FOLHAS INICIAIS
%===========================================
\newcommand{\folha}{
    \thispagestyle{empty}
    \vspace*{8cm}
    \begin{center}
        {\bfseries\scshape\Huge Folha de Rosto \par}
    \end{center}
    \clearpage
}

\newcommand{\contra}{
    \thispagestyle{empty}
    \vspace*{8cm}
    \begin{center}
        {\bfseries\scshape\Huge Contra Capa \par}
    \end{center}
    \clearpage
}

%===========================================
% CONFIGURAÇÃO DE CABEÇALHOS
%===========================================
\makepagestyle{meuestilo}
\makeevenhead{meuestilo}{\footnotesize\textit{\leftmark}}{}{\footnotesize\thepage}
\makeoddhead{meuestilo}{\footnotesize\textit{\leftmark}}{}{\footnotesize\thepage}
\makeheadrule{meuestilo}{\textwidth}{0.3pt}

% Redefinir marcas para incluir prefixos
\renewcommand{\chaptermark}[1]{%
    \markboth{\thechapter.\quad#1}{}%
}

% Primeira página de capítulo: apenas número no canto superior direito
\aliaspagestyle{chapter}{plain}
\makeoddhead{plain}{}{}{\footnotesize\thepage}
\makeevenhead{plain}{}{}{\footnotesize\thepage}

%===========================================
% CONFIGURAÇÃO DO SUMÁRIO
%===========================================
\renewcommand{\printtoctitle}[1]{%
    \begin{center}
        \huge\bfseries\MakeUppercase{#1}\\[0.5cm]
        \rule{\textwidth}{0.3pt}
    \end{center}
}
\renewcommand{\aftertoctitle}{\par\nobreak\vskip 2cm}
\setlength{\cftbeforechapterskip}{0.5cm}
\renewcommand{\cftchapterfont}{\bfseries\MakeUppercase}
\renewcommand{\cftchapterpagefont}{\bfseries}
\setlength{\cftchapterindent}{0cm}
\setlength{\cftsectionindent}{2cm}

%===========================================
% INÍCIO DO DOCUMENTO
%===========================================
\begin{document}

% Numeração arábica desde o início (contagem invisível)
\pagenumbering{arabic}
\pagestyle{empty}

\folha
\contra

%===========================================
% ELEMENTOS PRÉ-TEXTUAIS
%===========================================
\capituloCentralizado{Agradecimento}
Lorem ipsum dolor sit amet, consectetur adipiscing elit.

\capituloCentralizado{Resumo}
Este é o resumo do trabalho.

\capituloCentralizado{Lista de Figuras}
Conteúdo da lista de figuras.

\capituloCentralizado{Lista de Tabelas}
Conteúdo da lista de tabelas.

\capituloCentralizado{Lista de Abreviaturas e Siglas}
Conteúdo das abreviaturas.

\capituloCentralizado{Lista de Símbolos}
Conteúdo dos símbolos.

%===========================================
% SUMÁRIO
%===========================================
\tableofcontents
\clearpage

%===========================================
% PARTE TEXTUAL
%===========================================
\chapterstyle{CapitulosNumerado}
\pagestyle{meuestilo}

\capituloCentralizado{Introdução}
Este é o conteúdo da introdução do trabalho.

\part{Preparação do Relatório}
\chapter{Desenvolvimento}
Conteúdo do desenvolvimento.

\section{Primeira Seção}
Texto da primeira seção.

\section{Segunda Seção}
Texto da segunda seção.

\part{Resultado do Relatório}
\chapter{Conclusão}
Conteúdo da conclusão.

%===========================================
% ELEMENTOS PÓS-TEXTUAIS
%===========================================
\capituloCentralizado{Referências}
AUTOR, A. \textit{Título do livro}. Editora, 2024.

%===========================================
% APÊNDICES
%===========================================
\tituloDivisorio{Apêndices}

\appendix
\renewcommand{\chaptermark}[1]{%
    \markboth{APÊNDICE~\thechapter\quad--\quad#1}{}%
}

\chapter{Título do Apêndice A}
Conteúdo do apêndice A.

\section{Subseção A.1}
Conteúdo da subseção.

\chapter{Título do Apêndice B}
Conteúdo do apêndice B.

%===========================================
% ANEXOS
%===========================================
\tituloDivisorio{Anexos}

% Reset do contador para anexos
\setcounter{chapter}{0}
\renewcommand{\thechapter}{\Alph{chapter}}
\renewcommand{\chaptername}{Anexo}
\renewcommand{\chaptermark}[1]{%
    \markboth{ANEXO~\thechapter\quad--\quad#1}{}%
}

\chapter{Título do Anexo A}
Conteúdo do anexo A (numeração reiniciada).

\chapter{Título do Anexo B}
Conteúdo do anexo B.

\end{document}